% P15-U3: posterior laws, prediction, sensitivity and interpretation.
\chapter{Bayesian Updating and Prediction}
\label{sec:p15-bayesian-updating-prediction}

Bayesian inference starts with a probability model for observations
conditional on an unknown quantity and a probability distribution for
that quantity before those observations are used. Conditioning their
joint law produces a posterior distribution. Integrating a future
observation law against that posterior produces a prediction. Neither
operation supplies the sampling model, prior, observation mechanism or
decision losses: those are explicit inputs.

The posterior and predictive construction below develops the probability
identities used by Kass and Raftery
\cite[secs.~3.1--3.2 and 6.1]{KassRaftery1995BayesFactors},
using the measure and kernel foundations already established in
\cref{sec:p15-observation-kernel}. The sensitivity and interpretation
sections supply exact examples for the distinctions in their
Sections~5 and 8.2. The constructions and proofs here are editorial
developments; they do not purport to reproduce that article's separate
application studies or numerical integration algorithms.

\section{The joint model and its posterior}
\label{sec:p15-bayes-joint-model}

\begin{hbdefinition}{def:p15-bayes-experiment}{A dominated Bayesian experiment}
Let \((\Theta,\mathcal T)\) be a nonempty measurable parameter space,
\(\Pi\) a probability measure on it, and \((\mathcal X,\mathcal A,\nu)\)
a sigma-finite observation measure space. A likelihood specification is
a nonnegative, finite, jointly measurable function
\((\theta,x)\mapsto f_\theta(x)\) such that
\(\int_{\mathcal X}f_\theta(x)\,\nu(dx)=1\) for every \(\theta\).
The sampling law at a fixed parameter is
\(P_\theta(B)=\int_B f_\theta\,d\nu\).
For a fixed observation \(x\), the function
\(\theta\mapsto f_\theta(x)\) is its likelihood.

The prior predictive density and its regular observation set are
\[
 m(x)=\int_\Theta f_\theta(x)\,\Pi(d\theta),
 \qquad G=\{x:0<m(x)<\infty\}.
\]
The joint law \(Q\) of the parameter and data is
\[
 Q(A\times B)=\int_A P_\theta(B)\,\Pi(d\theta)
                 \quad(A\in\mathcal T,\ B\in\mathcal A).
\]
The parameter coordinate is random under \(Q\); it is fixed under
each sampling law \(P_\theta\). These are different probability spaces.
The prior is an input describing this particular inference problem;
the likelihood is not itself a probability distribution over \(\Theta\).
\index{Bayesian inference!experiment}\index{likelihood!parameter function}
\index{prior predictive distribution}
\end{hbdefinition}

\begin{hbtheorem}{thm:p15-bayes-posterior-kernel}{proved=annex:p15.bayes-posterior-kernel}{Normalization and a complete posterior kernel}
The observation marginal of \(Q\) has density \(m\), integrates to one,
and assigns probability one to \(G\). The formula
\[
 \Pi_x(A)=
 \begin{cases}
 \displaystyle\frac{\int_A f_\theta(x)\,\Pi(d\theta)}{m(x)},&x\in G,\\
 \Pi(A),&x\notin G
 \end{cases}
\]
defines a probability kernel from observations to parameters satisfying
\[
 Q(A\times B)=\int_B\Pi_x(A)m(x)\,\nu(dx).
\]
For \(x\in G\), \(\Pi_x\ll\Pi\) and its derivative is
\(f_\theta(x)/m(x)\). For every nonnegative measurable \(h\),
\[
 \int h(\theta)\,\Pi_x(d\theta)
   =\frac{\int h(\theta)f_\theta(x)\,\Pi(d\theta)}{m(x)}.
\]
The same identity holds for signed \(h\) when its posterior absolute
integral is finite.
\end{hbtheorem}

Values assigned off \(G\) do not describe conditioning on an event
of positive prior predictive probability. Even within \(G\), a
continuous observation has a density, rather than a positive singleton
probability. The displayed kernel supplies a version relative to the
specified likelihoods. Arbitrary changes of densities on null sets
cannot establish a uniquely compelled posterior at an exceptional
point. The explicit null-observation examples in
\cref{sec:p15-conditioning-coordinate-examples} explain this issue.

\section{Odds, evidence and dependence}
\label{sec:p15-bayes-odds}

\begin{hbdefinition}{def:p15-bayes-event-odds}{Comparing an event with its complement}
Fix \(H\in\mathcal T\) with \(p=\Pi(H)\in(0,1)\). At \(x\), define
\[
 \ell_H(x)=\frac1p\int_H f_\theta(x)\,\Pi(d\theta),\qquad
 \ell_{H^c}(x)=\frac1{1-p}\int_{H^c}f_\theta(x)\,\Pi(d\theta).
\]
These are observation densities averaged over the conditional priors,
not maximized likelihoods. When both are finite and positive, let
\(B_H(x)=\ell_H(x)/\ell_{H^c}(x)\).
\end{hbdefinition}

\begin{hbtheorem}{thm:p15-bayes-event-odds}{proved=annex:p15.bayes-event-odds}{The update multiplier and its endpoint cases}
For \(x\in G\),
\[
 \Pi_x(H)=\frac{p\ell_H(x)}
                 {p\ell_H(x)+(1-p)\ell_{H^c}(x)}.
\]
If both likelihoods are positive, then
\[
 \frac{\Pi_x(H)}{1-\Pi_x(H)}
       =B_H(x)\frac p{1-p}.
\]
In this case \(\Pi_x(H)>p\) exactly when \(B_H(x)>1\).
If just \(\ell_H\) vanishes the posterior is zero; if just
\(\ell_{H^c}\) vanishes it is one. If both vanish, \(x\notin G\).
A prior-null event remains posterior-null on \(G\).
\end{hbtheorem}

\begin{hbtheorem}{thm:p15-bayes-flag-calculation}{proved=annex:p15.bayes-flag-calculation}{A flag, a second measurement and a duplicated report}
Consider a stipulated two-state inspection model. State \(H\) has
prior probability \(1/1000\). A positive flag \(+\) has probabilities
\(99/100\) under \(H\) and \(1/100\) under \(H^c\).
One positive flag has likelihood ratio \(99\), but
\[
 P(H\mid +)=\frac{11}{122}.
\]
Two flags conditionally independent given the state have
\[
 P(H\mid +,+)=\frac{363}{400}.
\]
If the second report is instead a copy of the first, observing
\((+,+)\) leaves the posterior at \(11/122\).

With the same one-flag likelihoods and prior \(p\in(0,1)\),
\[
                    P(H\mid +)=\frac{99p}{1+98p}.
\]
Over \(p\in[1/10000,1/100]\) this ranges from \(1/102\) to \(1/2\).
The numerical values describe the supplied model, not a measured
inspection system.
\end{hbtheorem}

An observation can increase the probability of a claim while leaving
it unlikely. That probability can be high without being one, and a
probability-one conclusion within a model is not a logical consequence
across all models. The distinction between logical probability and
theory systematization is developed in
\cref{sec:p15-systematic-power}.

\section{Sequential updating}
\label{sec:p15-bayes-sequential}

\begin{hbdefinition}{def:p15-bayes-sequential-model}{Conditional sampling densities}
For a finite sequence \(x_{1:n}=(x_1,\ldots,x_n)\), let
\(f_i(x_i\mid\theta,x_{1:i-1})\) be a jointly measurable
sampling density relative to a sigma-finite measure \(\nu_i\),
normalized for every parameter and preceding history. Write
\[
 L_n(\theta)=\prod_{i=1}^n f_i(x_i\mid\theta,x_{1:i-1}),
 \qquad m_n(x_{1:n})=\int L_n(\theta)\,\Pi(d\theta),
 \qquad L_0=m_0=1.
\]
Assume the likelihood factors are finite. A regular observed path
is one for which \(0<m_i(x_{1:i})<\infty\) at every prefix.
Define \(\Pi_i(d\theta)=L_i(\theta)\Pi(d\theta)/m_i\).
All statements about a realized path below use this regularity.
\end{hbdefinition}

\begin{hbtheorem}{thm:p15-bayes-sequential-update}{proved=annex:p15.bayes-sequential-update}{Batch updating equals successive conditioning}
On a regular path,
\[
 \begin{gathered}
 r_i(x_i\mid x_{1:i-1})
     =\int f_i(x_i\mid\theta,x_{1:i-1})\,\Pi_{i-1}(d\theta)
     =\frac{m_i}{m_{i-1}},\\
 \Pi_i(d\theta)=
 \frac{f_i(x_i\mid\theta,x_{1:i-1})\,\Pi_{i-1}(d\theta)}
      {r_i(x_i\mid x_{1:i-1})},\qquad
 m_n=\prod_{i=1}^n r_i .
 \end{gathered}
\]
Thus \(\log m_n=\sum_i\log r_i\). The factors are conditional
predictive probabilities or densities. Replacing them by unconditional
one-observation marginals is generally invalid.
For conditionally independent observations under a fixed common
sampling density, reordering the observations leaves the posterior
unchanged. For history-dependent mechanisms such invariance is not asserted.
\end{hbtheorem}

\section{Prediction integrates parameter uncertainty}
\label{sec:p15-bayes-prediction}

\begin{hbdefinition}{def:p15-bayes-future-kernel}{A specified future observation law}
Let \(R(\theta,x,\cdot)\) be a probability kernel from
\(\Theta\times\mathcal X\) to a future measurable space
\((\mathcal Y,\mathcal B)\). This is the future law conditional
on the parameter and already observed data. Conditional independence
of past and future given the parameter is the special case
\(R(\theta,x,\cdot)=R_\theta(\cdot)\).
Define
\[
                       \widehat P_x(B)
                         =\int R(\theta,x,B)\,\Pi_x(d\theta).
\]
\index{posterior predictive distribution}\index{parameter uncertainty}
\end{hbdefinition}

\begin{hbtheorem}{thm:p15-bayes-predictive-kernel}{proved=annex:p15.bayes-predictive-kernel}{Prediction, conditional means and total variance}
The formula above is a probability kernel and is the conditional
future law in the joint model constructed from \(Q\) and \(R\).
For a real future variable \(Y\) with finite posterior predictive
second moment, put
\[
 a(\theta,x)=E(Y\mid\theta,x),\qquad
 b(\theta,x)=\operatorname{Var}(Y\mid\theta,x).
\]
These conditional moments are needed only almost everywhere under
\(\Pi_x\); assign zero on any exceptional set.
Then
\[
 \begin{split}
 E(Y\mid x)&=\int a(\theta,x)\,\Pi_x(d\theta),\\
 \operatorname{Var}(Y\mid x)
  &=\int b(\theta,x)\,\Pi_x(d\theta)
       +\operatorname{Var}_{\Pi_x}(a(\theta,x)).
 \end{split}
\]
For two future variables conditionally independent given \((\theta,x)\),
their predictive covariance equals the posterior covariance of their
conditional means whenever their second moments are finite.
\end{hbtheorem}

\begin{hbtheorem}{thm:p15-bayes-beta-calculation}{proved=annex:p15.bayes-beta-calculation}{Bernoulli learning and several future draws}
Let \(\theta\sim\operatorname{beta}(\alpha,\beta)\), with
\(\alpha,\beta>0\), and let \(X_1,\ldots,X_n\) be conditionally
independent Bernoulli variables of success probability \(\theta\).
Allow \(n=0\). For \(k=\sum_i x_i\), set
\(a=\alpha+k,\ b=\beta+n-k,\ t=a+b\). Then
\[
 \Pi_{x_{1:n}}=\operatorname{beta}(a,b),\qquad
 m_n(x_{1:n})=\frac{B(a,b)}{B(\alpha,\beta)}.
\]
For the count \(K_n\), rather than a particular binary word, the
predictive probability is \(\binom nk B(a,b)/B(\alpha,\beta)\).
Both observations give the same posterior when the sampling model
records the same fixed sample size.

For \(r\geq1\) further draws, with future count \(J_r\),
\[
 \begin{gathered}
 P(J_r=j\mid x_{1:n})
    =\binom rj\frac{B(a+j,b+r-j)}{B(a,b)},\quad 0\leq j\leq r,\\
 E(J_r\mid x_{1:n})=\frac{ra}{t},\qquad
 \operatorname{Var}(J_r\mid x_{1:n})
     =\frac{rab(t+r)}{t^2(t+1)}.
 \end{gathered}
\]
Under the uniform prior and \(k=8,n=10\), the posterior is
\(\operatorname{beta}(9,3)\), with mean \(3/4\), variance \(3/208\),
and mode \(4/5\). The probability that two new draws both succeed
is \(15/26\), rather than the plug-in value \(9/16\).
The earlier urn construction in
\cref{thm:p15-polya-posterior-prediction} supplies the distinct
reinforcement interpretation of this same beta mixture.
\end{hbtheorem}

\begin{hbtheorem}{thm:p15-bayes-normal-calculation}{proved=annex:p15.bayes-normal-calculation}{Normal learning with known observation variance}
Let \(n\geq1,\ \sigma^2>0,\ \tau^2>0\).
Given \(\theta\), take independent
\(X_i\sim\operatorname{normal}(\theta,\sigma^2)\), and put
\(\theta\sim\operatorname{normal}(\mu_0,\tau^2)\) before observing them.
Write
\[
 \bar x=\frac1n\sum_i x_i,\quad
 S=\sum_i(x_i-\bar x)^2,\quad s^2=\frac{\sigma^2}{n},\quad
 v=(\tau^{-2}+s^{-2})^{-1},\quad
 u=v(\mu_0\tau^{-2}+\bar x s^{-2}).
\]
The posterior is \(\operatorname{normal}(u,v)\). With
\(\phi_V(z)=(2\pi V)^{-1/2}\exp(-z^2/(2V))\) for \(V>0\),
the joint observation density is
\[
 m_n(x_{1:n})
 =\frac{(2\pi\sigma^2)^{-(n-1)/2}}{\sqrt n}
       \exp\!\left(-\frac{S}{2\sigma^2}\right)
       \phi_{\tau^2+s^2}(\bar x-\mu_0).
\]
One future observation has law
\(\operatorname{normal}(u,\sigma^2+v)\); the average of \(r\geq1\)
new conditionally independent observations has law
\(\operatorname{normal}(u,\sigma^2/r+v)\).
Distinct future observations have posterior predictive covariance \(v\).

For \(\mu_0=0,\tau^2=1,\sigma^2=4,n=4,\bar x=2\), the posterior
is \(\operatorname{normal}(1,1/2)\), and the one-observation
predictive law is \(\operatorname{normal}(1,9/2)\).
Neither is the plug-in law obtained by setting \(\theta=\bar x\).
\end{hbtheorem}

\section{Reparameterization and prior sensitivity}
\label{sec:p15-bayes-prior-sensitivity}

\begin{hbtheorem}{thm:p15-bayes-transformation}{proved=annex:p15.bayes-transformation}{Transforming an inferred quantity and changing its prior}
For any measurable map \(g:(\Theta,\mathcal T)\to(S,\mathcal S)\),
the posterior law of \(g(\theta)\)
is the pushforward \(g_\#\Pi_x\), that is
\((g_\#\Pi_x)(A)=\Pi_x(g^{-1}(A))\).
If \(g\) is a measurable bijection with measurable inverse and both
the prior and sampling family are transported through it, the posterior
and predictive probabilities are unchanged by this reparameterization.

A uniform prior for \(\theta\in(0,1)\) induces for
\(\eta=\theta^2\) the density \(1/(2\sqrt\eta)\).
A uniform prior for \(\eta\) instead induces the density \(2\theta\)
for \(\theta\). After one Bernoulli success, the respective posterior
means of \(\theta\) are \(2/3\) and \(3/4\), and those of \(\eta\)
are \(1/2\) and \(3/5\). Thus specifying uniformity separately in
two coordinates generally changes the inference.
\end{hbtheorem}

\begin{hbtheorem}{thm:p15-bayes-reweighting}{proved=annex:p15.bayes-reweighting}{Exact sensitivity by reweighting and contamination}
Let \(w:\Theta\to[0,\infty)\) be measurable with
\(0<\Pi w<\infty\), where \(\Pi w=\int w\,d\Pi\), and set
\(\Pi^w(d\theta)=w(\theta)\Pi(d\theta)/\Pi w\).
At a regular observation with \(0<\Pi_xw<\infty\),
\[
 \Pi_x^w(d\theta)=\frac{w(\theta)\Pi_x(d\theta)}{\Pi_xw},
 \qquad \frac{m^w(x)}{m(x)}=\frac{\Pi_xw}{\Pi w}.
\]
If \(0<a\leq w\leq b<\infty\) almost everywhere under \(\Pi\),
then for any \(H\in\mathcal T\) and \(p_x=\Pi_x(H)\),
\[
 \frac{ap_x}{ap_x+b(1-p_x)}
 \leq\Pi_x^w(H)\leq
 \frac{bp_x}{bp_x+a(1-p_x)}.
\]
Both bounds can be attained by weights constant on \(H\) and \(H^c\).

Alternatively, for \(\Pi^\varepsilon=(1-\varepsilon)\Pi+\varepsilon R\),
\(0<\varepsilon<1\), suppose \(R\) is a probability prior with
finite positive predictive density \(m_R(x)\). Then
\[
 \Pi_x^\varepsilon=(1-\lambda_x)\Pi_x+\lambda_x R_x,\qquad
 \lambda_x=\frac{\varepsilon m_R(x)}
                  {(1-\varepsilon)m(x)+\varepsilon m_R(x)}.
\]
Small prior contamination need not remain small after conditioning.
\end{hbtheorem}

\section{Credible sets and repeated sampling}
\label{sec:p15-bayes-credible-sets}

\begin{hbdefinition}{def:p15-bayes-credible-set}{Posterior content and coverage}
Fix \(0<\alpha<1\). A measurable set-valued procedure \(C\)
means that \(\{(\theta,x):\theta\in C(x)\}\) is product measurable.
It has posterior content at least \(1-\alpha\) when
\(\Pi_x(C(x))\geq1-\alpha\) for almost every observation under \(m\nu\).
Its frequentist coverage at fixed \(\theta\) is
\[
                     c(\theta)=P_\theta\{\theta\in C(X)\}.
\]
A confidence guarantee at level \(1-\alpha\) requires
\(c(\theta)\geq1-\alpha\) for every parameter in its stated scope.
Posterior content concerns the conditional parameter law after
observing data; coverage concerns the procedure under repeated data
at a fixed parameter.
\index{credible set}\index{confidence set!comparison with credibility}
\end{hbdefinition}

\begin{hbtheorem}{thm:p15-bayes-coverage}{proved=annex:p15.bayes-coverage}{Prior-average calibration and failure of uniform coverage}
For any such \(C\),
\[
 \int_\Theta c(\theta)\,\Pi(d\theta)
   =\int_{\mathcal X}\Pi_x(C(x))m(x)\,\nu(dx).
\]
Thus posterior content at least \(1-\alpha\) implies prior-average
coverage at least \(1-\alpha\); exact content implies exact average.

This does not imply a uniform fixed-parameter coverage bound.
For one Bernoulli observation and a uniform prior, the equal-tail
posterior intervals of content \(0.95\) are
\[
 \begin{split}
 C(1)&=[\sqrt{0.025},\,\sqrt{0.975}],\\
 C(0)&=[1-\sqrt{0.975},\,1-\sqrt{0.025}].
 \end{split}
\]
At the fixed value \(\theta=0.99\), both intervals miss the parameter,
so coverage is zero despite exact posterior content \(0.95\)
for either observation.
\end{hbtheorem}

For the normal posterior above, let \(\Phi\) denote the standard
normal CDF and \(z_{1-\alpha/2}\) its unique
\((1-\alpha/2)\)-quantile. Then
\([u-z_{1-\alpha/2}\sqrt v,u+z_{1-\alpha/2}\sqrt v]\)
has posterior content \(1-\alpha\). The sampling coverage of this
particular procedure must be calculated under its fixed-parameter law.
The posterior normalization alone does not calculate it.

\section{Decisions and predictive scores}
\label{sec:p15-bayes-decisions}

\begin{hbdefinition}{def:p15-bayes-loss}{Loss and posterior risk}
For a finite nonempty action set \(\mathcal D\), let
\(L(d,\theta)\) be a bounded nonnegative measurable loss.
The posterior risk is
\(\rho_x(d)=\int L(d,\theta)\,\Pi_x(d\theta)\).
A Bayes action minimizes \(\rho_x\). An ordering of the finite
action set resolves ties and makes the resulting rule measurable.
The choice of a loss describes the consequences to be assessed;
conditioning alone does not determine that choice.
\end{hbdefinition}

\begin{hbtheorem}{thm:p15-bayes-decision-rule}{proved=annex:p15.bayes-decision-rule}{Optimal actions and expected value of the observation}
A posterior risk minimizer also minimizes the prior integrated
expected loss among all measurable rules with values in \(\mathcal D\).
Furthermore,
\[
 E\!\left[\min_{d\in\mathcal D}\rho_X(d)\right]
       \leq \min_{d\in\mathcal D}\int L(d,\theta)\,\Pi(d\theta).
\]
For actions \emph{declare \(H\)} and \emph{declare \(H^c\)}, with
loss \(c_{\rm FP}>0\) for the first action when \(H^c\) is true,
loss \(c_{\rm FN}>0\) for the second when \(H\) is true, and zero
otherwise, declaring \(H\) is optimal exactly when
\[
            \Pi_x(H)\geq\frac{c_{\rm FP}}{c_{\rm FP}+c_{\rm FN}}.
\]
For a real parameter and squared-error loss, if
\(\int\theta^2\,\Pi_x(d\theta)<\infty\), the unique minimizing
real action is its posterior mean.
\end{hbtheorem}

In the flag example, equal losses favor declaring \(H^c\) after one
positive. If \(c_{\rm FN}=20\) and \(c_{\rm FP}=1\), the threshold
is \(1/21\), below \(11/122\), and the optimal declaration changes.
The posterior evidence has not changed; the consequences have.
Preference conditions supporting expected-utility evaluation belong
to \cref{sec:p17-finite-expected-utility}.

\begin{hbtheorem}{thm:p15-bayes-log-score}{proved=annex:p15.bayes-log-score}{A finite predictive probability is optimal for its logarithmic score}
On a finite outcome set let \(p\) be the specified predictive mass
function. A report \(q\) is any probability mass function on that set.
Define its expected logarithmic loss as
\[
              \mathcal L_p(q)=-\sum_{y:p(y)>0}p(y)\log q(y),
\]
with value \(+\infty\) if \(q(y)=0<p(y)\).
Then
\[
 \mathcal L_p(q)-\mathcal L_p(p)
       =\sum_{y:p(y)>0}p(y)\log\frac{p(y)}{q(y)}\geq0,
\]
and equality holds exactly when \(q=p\).
For finite observation sequences with finite positive predictive
probabilities at every observed prefix, the cumulative conditional
logarithmic score is \(-\log m_n(x_{1:n})\).
This optimality is relative to the stipulated predictive law \(p\);
it does not prove that \(p\) describes the physical data source.
\end{hbtheorem}

\section{Stopping, selection and what was observed}
\label{sec:p15-bayes-stopping-selection}

\begin{hbtheorem}{thm:p15-bayes-stopping}{proved=annex:p15.bayes-stopping}{When an observation mechanism cancels, and when it does not}
Suppose an observed finite record \(z\), including its length, has
likelihood
\[
              P_\theta(Z=z)=c(z)L_z(\theta),
                  \qquad 0<c(z)<\infty ,
\]
where \(L_z\) is nonnegative and measurable and \(c(z)\) does not
depend on \(\theta\). When the posterior
normalizer is finite and positive, the posterior is proportional
to \(L_z(\theta)\Pi(d\theta)\). Such a common factor cancels.
If the recording probability is instead \(c_\theta(z)\),
that factor must be included.

For independent Bernoulli trials stopped at their first success,
a beta prior with positive shapes gives, after \(N=n\),
\[
                   \theta\mid N=n
                     \sim\operatorname{beta}(\alpha+1,\beta+n-1).
\]
This equals the posterior from the same fixed-length word of
\(n-1\) failures and one success. It also equals the posterior
from a fixed-size count of one success in \(n\) trials, whose
likelihood has the extra parameter-independent factor \(n\).

For a different experiment that samples a Bernoulli outcome
\emph{conditional on it being positive}, its reported value is
always one for every \(\theta\in(0,1)\) and carries constant
likelihood. With a uniform prior it leaves the prior unchanged.
Observing a positive outcome from one predesignated, unselected
trial instead gives the \(\operatorname{beta}(2,1)\) posterior.
\end{hbtheorem}

To use the cancellation result, specify the experiment and all recorded
information. Recording that a predesignated trial was selected is
different from receiving a value drawn from a selected-only population.
The theorem concerns the posterior for the stated record; it supplies
no unrestricted frequentist error guarantee under optional stopping.
It also gives no permission to reuse data as a prior and then count
them again in the likelihood.

\section{A procedure for an inference problem}
\label{sec:p15-bayes-reference-procedure}

\begin{hbmethod}{def:p15-bayes-reference-procedure}{From assumptions to a posterior report}
State the unknown quantity, observation mechanism, sampling law and
prior. Identify whether the record is a complete sequence, a count,
a stopped experiment or a selected observation. Form the joint
likelihood with its dependence structure. Check that the posterior
normalizing integral is positive and finite.

Compute the posterior and the requested transformed quantities or
future law. Use a joint predictive law for several future observations,
rather than silently replacing a shared parameter by its estimate.
Report posterior probability, a credible set, a prediction or a
loss-based action with its corresponding meaning.

Vary material prior weights and sampling assumptions using the
reweighting identity or direct recalculation. Keep likelihood ratios,
posterior probabilities, confidence coverage and decision thresholds
explicitly distinguished. An impossible observation, an improper
normalizer or an omitted observation mechanism requires a changed
model before these formulas can be used.
\end{hbmethod}

For comparing several fully specified models and averaging their
predictions, continue with \cref{sec:p15-bayesian-model-comparison}.
For controlled sampling error, randomized tests and confidence-set
construction, use \cref{sec:p15-hypothesis-tests-confidence}.
Those methods require their own inputs and establish their stated
conclusions; the posterior calculation is not a substitute for their
sampling guarantees.
